\section{BR-123-OPPG2: every containing Borel pair is opposite in type \texorpdfstring{$G_2$}{G2}}

For the original question in cote 123, physical page 22/footer 21,
the answer to ``all choices'' is affirmative in type $G_2$.
We prove this using both fundamental representations, keeping the
restriction from the ambient orthogonal group explicit.

\subsection{The seven-dimensional cyclic model}
Let $\alpha$ be the short simple root and $\beta$ the long one,
normalized by
\[
 (\alpha,\alpha)=2,\qquad(\beta,\beta)=6,
                  \qquad(\alpha,\beta)=-3.
\]
The fundamental weights are $\omega_1=2\alpha+\beta$ and
$\omega_2=3\alpha+2\beta=\theta$.  Let $V=L(\omega_1)$.
The Weyl dimension formula, Etingof (n.d.), Proposition 28.1,
p.~125, gives
\[
 \dim V=\frac21\frac33\frac54\frac75\frac96\frac{12}9=7.
 \tag{G2.1}
\]
Here the positive roots in order are
$\alpha,\beta,\alpha+\beta,2\alpha+\beta,3\alpha+\beta,
3\alpha+2\beta$, and $\rho=5\alpha+3\beta$.
The Weyl orbit of $\omega_1$ consists of the six short roots,
each of multiplicity one.  The single remaining weight must be
Weyl-fixed, hence zero.  Thus all seven weights have multiplicity
one.  Since $w_0=-1$, the module is self-dual by Etingof,
Proposition 33.5, p.~148.  Its nondegenerate invariant bilinear
form is symmetric: its transpose differs by a scalar by Schur's
lemma, so it is symmetric or alternating, and a nondegenerate
alternating form cannot have odd dimension.  Denote this form by $B$.

Choose a principal $\mathfrak{sl}_2$ triple with raising element
$N=e_\alpha+e_\beta$ and semisimple element $2\rho^\vee$;
the construction is Etingof, \S33.3, p.~149.  The weights of
$2\rho^\vee$ on $V$ are $-6,-4,-2,0,2,4,6$.
The highest-weight-six $\mathfrak{sl}_2$ summand already has
dimension seven, so the restriction is that single irreducible
module.  Consequently a weight basis $e_0,\ldots,e_6$ can be
chosen with
\[
 N=\sum_{i=0}^{5}E_{i+1,i},\qquad
 B(e_i,e_j)=(-1)^i\delta_{i+j,6}.
 \tag{G2.2}
\]
Indeed the first normalization uses the nonzero successive raising
coefficients, and invariance of $B$ then gives alternating successive
anti-diagonal entries; scaling $B$ sets the first entry to one.
The weight list, in this basis order, is
\[
 -(2\alpha+\beta),\quad-(\alpha+\beta),\quad-\alpha,
 \quad0,\quad\alpha,\quad\alpha+\beta,\quad2\alpha+\beta.
 \tag{G2.3}
\]

The representation embeds $\mathfrak g_2$ into
$\mathfrak{so}(V,B)$, since its kernel is an ideal and the
representation is nontrivial.  The grade $-5$ part of the latter
Lie algebra for $\rho^\vee$ is one-dimensional, spanned by
\[
                         Y=E_{0,5}+E_{1,6}.
 \tag{G2.4}
\]
To check this directly, a matrix entry of grade $-5$ has row
index five less than its column index, so only $(0,5)$ and
$(1,6)$ occur; the skew-adjointness equation for (G2.2) equates
their coefficients.  The nonzero negative-highest-root space
of $\mathfrak g_2$ has precisely this grade and therefore maps
onto this line.  In particular $X=N+Y$ belongs to
$\mathfrak g_2$ and has all simple and negative-highest-root
coefficients nonzero.  Kostant (2010), Theorems 1.2--1.3,
p.~6, identifies its centralizing torus $T_1$ as a torus in
apposition to the original weight torus $T_0$.

Put $\zeta=e^{2\pi i/6}$, choose $a$ with $a^6=2$, and set
$\lambda_j=a\zeta^j$ for $j\in\mathbf Z/6$.  The seven
eigenvectors are
\[
 v_j=e_0+\sum_{i=1}^{5}\lambda_j^{6-i}e_i+e_6,
       \quad Xv_j=\lambda_jv_j,
       \qquad w=e_0-e_6,\quad Xw=0.
 \tag{G2.5}
\]
These identities follow by multiplication using (G2.2)--(G2.4):
the row-one equation is $2=\lambda_j^6$, the interior equations
shift a power, and the boundary equations are immediate.
The distinct eigenvalues prove that these vectors form a basis.
The lines $\mathbf C v_j$ are exactly the six nonzero weight
lines of $T_1$, and $\mathbf Cw$ is its zero-weight line.

The principal element $c=\rho^\vee(\zeta)$ satisfies
$cXc^{-1}=\zeta X$, so $c$ permutes these six lines cyclically
by $j\mapsto j-1$.  On $T_1$ it acts as a Coxeter element.
In the rank-two $G_2$ root system a Coxeter element rotates the
short-root hexagon by $60$ degrees in one orientation: the
simple-reflection product has order six and trace one on the
root plane, as follows immediately from the above inner products.
Thus consecutive indices $j,j+1$ correspond to adjacent short
roots in the hexagon.  Their sum is a long root.

\subsection{Both fundamental highest-line pairings}
Under the canonical equivariant identification
$\mathfrak{so}(V,B)\simeq\bigwedge^2V$, the adjoint
representation of $G_2$ is an irreducible subrepresentation
of highest weight $\omega_2=\theta$.  Each long-root weight
occurs just once in $\bigwedge^2V$, from the unique unordered
pair of short roots whose sum is that long root.  In the
$T_0$ weight basis these six lines are
\[
 \mathbf C(e_i\wedge e_k),\qquad
 (i,k)\in\{(0,1),(0,2),(1,4),(2,5),(4,6),(5,6)\}.
 \tag{G2.6}
\]
This list is verified by adding the weights in (G2.3).
In the $T_1$ weight basis they are exactly
\[
                         \mathbf C(v_j\wedge v_{j+1}),
                         \qquad j\in\mathbf Z/6.
 \tag{G2.7}
\]
Indeed the six adjacent sums are the six distinct long roots,
and each has multiplicity one, so its wedge line is the
corresponding line in the adjoint subrepresentation.

The form $B$ induces
\[
 B_2(u\wedge u',v\wedge v')
        =B(u,v)B(u',v')-B(u,v')B(u',v)
 \tag{G2.8}
\]
on $\bigwedge^2V$.  Its restriction to the adjoint
subrepresentation is nondegenerate: it pairs the lines
$e_5\wedge e_6$ and $e_0\wedge e_1$ nontrivially,
so it is a nonzero invariant form on an irreducible module.

First, all entries $B(e_i,v_j)$ with
$i\in\{0,1,2,4,5,6\}$ are nonzero by (G2.2) and (G2.5).
Consequently every pair of highest lines for the two tori in
the $\omega_1$ representation has nonzero pairing.

For the $\omega_2$ representation, write
\[
 (r_0,r_1,r_2,r_3,r_4,r_5,r_6)=(0,5,4,3,2,1,0).
\]
The coefficient of $e_i$ in $v_j$ is $\lambda_j^{r_i}$.
For $(i,k)$ in (G2.6), (G2.8) is a nonzero sign times
\[
 \det\begin{pmatrix}
  \lambda_j^{r_{6-i}}&\lambda_{j+1}^{r_{6-i}}\\
  \lambda_j^{r_{6-k}}&\lambda_{j+1}^{r_{6-k}}
 \end{pmatrix}
  =\lambda_j^{r_{6-i}+r_{6-k}}
        (\zeta^{r_{6-k}}-\zeta^{r_{6-i}}).
 \tag{G2.9}
\]
The residues $r_s$ are distinct except for $s=0,6$,
and none of the pairs in (G2.6) is $(0,6)$.
Thus (G2.9) never vanishes.  Every pair of highest lines
in the second fundamental representation also has nonzero
pairing.

We finish by proving why these two tests suffice.  Fix a
Borel $B_+$ and highest vectors $z_1,z_2$ in its two
fundamental representations.  Their invariant forms pair
opposite weights nondegenerately.  For another Borel
$gB_+g^{-1}$ write $g=b_1n_ub_2$ by the Bruhat
decomposition, with $b_1,b_2\in B_+$ and $u\in W$.
Since $B_+$ fixes each highest line, invariance gives
\[
             (z_i,gz_i)=\text{nonzero scalar}\,(z_i,n_uz_i).
\]
The last pairing is nonzero exactly when
$u\omega_i=-\omega_i$.  Therefore nonzero pairings for
both $i=1,2$ force $u=-1=w_0$, because the fundamental
weights form a basis of the root plane.  This is exactly
the opposite-Borel Bruhat cell.  Conversely that cell
has the stated nonzero pairings.  This argument uses
Milne (2017), Theorem 21.80, p.~450, for Bruhat
decomposition and Theorem 21.84, p.~451, for opposite Borels.

For any $B_0\supset T_0$ and $B_1\supset T_1$, their
fundamental highest lines are among the lines just tested.
Thus $B_0,B_1$ are opposite.  CC-APG01 transfers this result
from the constructed ordered pair to every apposition pair
in type $G_2$.  This proves the assertion. $\square$

\paragraph{Primary representation-theoretic reference.}
Etingof, P. (n.d.). \emph{18.745: Lie groups and Lie algebras, I}
[Course notes, 223-page author PDF]. Massachusetts Institute
of Technology. \url{https://math.mit.edu/~etingof/lnlg.pdf}.
Proposition 28.1, p.~125; Proposition 33.5, p.~148;
\S33.3, p.~149.  The finite dimension product and the
matrix and minor calculations in this section are carried
out here; no result about apposition in this source is assumed.
